% part: intuitionistic-logic
% chapter: introduction
% section: bhk-interpretation

\documentclass[../../../include/open-logic-section]{subfiles}

\begin{document}

\olfileid{int}{int}{bhk}

\olsection{Brouwer--Heyting--Kolmogorov 解释}

\begin{editorial}
  使用 BHK 解释证明直觉主义命题的有效性令人困惑；需要对其作更好的解释。
\end{editorial}

对于直觉主义联结词，存在一种非形式的构造性解释，通常称为 Brouwer--Heyting--Kolmogorov解释。它使用了“构造”这一概念，你可以将其理解为一种构造性证明。（我们在 BHK 解释中不使用“证明”一词，以免与形式推导系统中的推导概念相混淆。）基于这一直观概念，BHK 解释阐明了直觉主义联结词的意义。

\begin{enumerate}
\item 我们假设已知何为原子语句的构造。
\item $!A_1 \land !A_2$ 的构造是一个序对 $\tuple{M_1, M_2}$，其中 $M_1$ 是 $!A_1$ 的构造，而 $M_2$ 是 $!A_2$ 的构造。
\item $!A_1 \lor !A_2$ 的构造是一个序对 $\tuple{s, M}$，其中 $s$ 为 $1$ 且 $M$ 是 $!A_1$ 的构造，或者 $s$ 为 $2$ 且 $M$ 是 $!A_2$ 的构造。
\item $!A \lif !B$ 的构造是一个函数，它能将 $!A$ 的构造转换为 $!B$ 的构造。
\item 不存在 $\lfalse$（假）的构造。
\item $\lnot !A$ 定义为 $!A \lif \lfalse$ 的同义词。也就是说，$\lnot !A$ 的构造是一个函数，它将 $!A$ 的构造转换为 $\lfalse$ 的构造。
\end{enumerate}

\begin{ex}
以 $\lnot \lfalse$ 为例。它的构造是一个函数：给定任何 $\lfalse$ 的构造作为输入，就能输出一个 $\lfalse$ 的构造。显然，恒等函数 $\Id{}$ 就是这样的一个构造：给定一个 $\lfalse$ 的构造 $M$，$\Id{}(M) = M$ 就给出了一个 $\lfalse$ 的构造。
\end{ex}

一般地说，$\lnot !A$ 意味着“$!A$ 的构造是不可能的”。

\begin{ex}
让我们为任意命题 $!A$ 给出 $!A \lif \lnot \lnot !A$ 的一个构造，即 $!A \lif ((!A \lif \lfalse) \lif \lfalse)$。这个构造应当是一个函数 $f$，它接受一个 $!A$ 的构造 $M$，并返回一个 $(!A \lif \lfalse) \lif \lfalse$ 的构造 $f(M)$。以下是 $f$ 给出 $(!A \lif \lfalse) \lif \lfalse$ 的构造的方式：我们需要定义一个函数 $g$，当输入一个 $!A \lif \lfalse$ 的构造 $h$ 时，它输出一个 $\lfalse$ 的构造。我们可以如下定义 $g$：将输入 $h$ 应用于我们之前得到的那个 $!A$ 的构造 $M$。由于 $h$ 的输出 $h(M)$ 是 $\lfalse$ 的一个构造，因此若 $M$ 是 $!A$ 的构造，$f(M)(h) = h(M)$ 便是 $\lfalse$ 的一个构造。
\end{ex}

\begin{ex}
让我们为 $\lnot(!A \land \lnot !A)$（即 $(!A
\land (!A \lif \lfalse)) \lif \lfalse$）给出一个构造。这是一个函数 $f$，给定一个 $!A \land (!A \lif \lfalse)$ 的构造 $M$ 作为输入，产生一个 $\lfalse$ 的构造。一个合取式 $!B_1 \land !B_2$ 的构造是一个序对 $\tuple{N_1, N_2}$，其中 $N_1$ 是 $!B_1$ 的构造，$N_2$ 是 $!B_2$ 的构造。我们可以定义函数 $p_1$ 和 $p_2$，分别从一个 $!B_1 \land !B_2$ 的构造中提取出 $!B_1$ 和 $!B_2$ 的构造：
\begin{align*}
  p_1(\tuple{N_1, N_2}) &= N_1\\
  p_2(\tuple{N_1, N_2}) &= N_2
\end{align*}
以下是 $f$ 的做法：首先，将 $p_1$ 应用于其输入 $M$，由此得到 $!A$ 的一个构造。然后，将 $p_2$ 应用于 $M$，得到 $!A \lif \lfalse$ 的一个构造。而这样的构造本身又是一个函数 $p_2(M)$，若输入 $!A$ 的一个构造，它将输出 $\lfalse$ 的一个构造。换句话说，如果我们将 $p_2(M)$ 应用于 $p_1(M)$，就得到了一个 $\lfalse$ 的构造。因此，我们可以定义 $f(M) = p_2(M)(p_1(M))$。
\end{ex}

\begin{ex}
让我们给出 $((!A \land !B) \lif !C) \lif (!A \lif
(!B \lif !C))$ 的一个构造，即一个函数 $f$，它能把 $(!A \land !B) \lif !C$ 的一个构造~$g$ 转换为 $(!A \lif (!B \lif
!C))$ 的一个构造。构造 $g$ 本身是一个函数（从 $!A \land !B$ 的构造到 $!C$ 的构造）。而输出 $f(g)$ 是一个函数~$h_g$，它把 $!A$ 的构造映射为从 $!B$ 的构造到 $!C$ 的构造的函数。

确实，这有点令人困惑。我们必须构造一个特定的函数 $h_g$，它将是 $f$ 在输入~$g$ 时的输出。$h_g$ 的输入是 $!A$ 的一个构造~$M$。$h_g(M)$ 的输出应当是一个函数~$k_M$，它从 $!B$ 的构造~$N$ 映射到 $!C$ 的构造。令 $k_{g,M}(N) = g(\tuple{M, N})$。记住 $\tuple{M, N}$ 是 $!A \land !B$ 的一个构造。所以 $k_{g,M}$ 是 $!B \lif !C$ 的一个构造：它将 $!B$ 的构造~$N$ 映射到 $!C$ 的构造。现在令 $h_g(M) = k_{g,M}$。这是一个把 $!A$ 的构造~$M$ 映射为 $!B \lif !C$ 的构造 $k_{g,M}$ 的函数。现在令 $f(g) = h_g$。这是一个将 $(!A \land !B) \lif !C$ 的构造 $g$ 映射到 $!A \lif (!B \lif !C)$ 的构造的函数。呼！{}
\end{ex}

语句 $!A \lor \lnot !A$ 被称为排中律。我们可以对某些特定的 $!A$ 证明它（例如 $\lfalse \lor \lnot \lfalse$），但无法一般地证明。这是因为直觉主义析取要求给出其中一个析取支的构造，但存在一些目前既不能被证明也不能被反驳的语句（比如哥德巴赫猜想）。然而，排中律也是无法反驳的：亦即，$\lnot \lnot (!A \lor \lnot !A)$ 成立。

\begin{ex}
为了证明 $\lnot \lnot (!A \lor \lnot !A)$，我们需要一个函数~$f$，它把 $\lnot(!A \lor \lnot !A)$ 的一个构造，即 $(!A
\lor (!A \lif \lfalse)) \lif \lfalse$，转换成一个~$\lfalse$ 的构造。换句话说，我们需要一个函数 $f$，使得如果 $g$ 是 $\lnot(!A \lor \lnot !A)$ 的一个构造，那么 $f(g)$ 就是~$\lfalse$ 的一个构造。

假设 $g$ 是 $\lnot(!A \lor \lnot !A)$ 的一个构造，即一个函数，它把 $!A \lor \lnot !A$ 的构造转换为~$\lfalse$ 的构造。$!A \lor \lnot !A$ 的一个构造是一个序对 $\tuple{s, M}$，其中要么 $s = 1$ 且 $M$ 是 $!A$ 的一个构造，要么 $s = 2$ 且 $M$ 是 $\lnot !A$ 的一个构造。令 $h_1$ 为将 $!A$ 的一个构造~$M_1$ 映射到 $!A \lor \lnot !A$ 的一个构造的函数：它将 $M_1$ 映射到 $\tuple{1, M_1}$。再令 $h_2$ 为将 $\lnot !A$ 的一个构造~$M_2$ 映射到 $!A \lor \lnot !A$ 的一个构造的函数：它将 $M_2$ 映射到 $\tuple{2, M_2}$。

令 $k$ 为 $\comp{h_1}{g}$：它是一个函数，若给定 $!A$ 的一个构造，就返回一个~$\lfalse$ 的构造，也就是说，它是 $!A \lif \lfalse$ 即 $\lnot !A$ 的一个构造。现在令 $l$ 为 $\comp{h_2}{g}$。它是一个函数，若给定 $\lnot !A$ 的一个构造，就给出一个~$\lfalse$ 的构造。由于 $k$ 是 $\lnot !A$ 的一个构造，所以 $l(k)$ 是~$\lfalse$ 的一个构造。

总而言之，我们所做的就是描述了如何将 $\lnot(!A \lor \lnot !A)$ 的一个构造~$g$ 转换为一个~$\lfalse$ 的构造，亦即，将 $\lnot(!A \lor \lnot !A)$ 的构造~$g$ 映射到~$\lfalse$ 的构造 $l(k)$ 的函数 $f$，就是 $\lnot\lnot(!A \lor \lnot !A)$ 的一个构造。
\end{ex}

正如你所见，使用 BHK 解释来证明公式的直觉主义有效性很快就会变得繁琐且令人困惑。幸运的是，对于直觉主义逻辑存在着更好的推导系统，以及更精确的语义解释。
\end{document}